\input{preamble.tex}

\title{\huge{Hand in}}
\author{\LARGE{Thobias Høivik}}
\date{}

\begin{document}
\maketitle

\newpage
\begin{prob*}[1]
    Let $X = [-1,1] \times \{-1,+1\} \subseteq \mathbb R^2$ have 
    the usual metric topology and define a relation 
    $\sim$ on $X$ saying that $(a,-1) \sim (a,1)$ for all 
    $a \in [-1,0) \cup (0,1]$. Let $Y = X /\sim$ have 
    the quotient topology 

    \begin{itemize}
        \item Check that $\sim$ is an equivalence relation 
        \item Is $Y$ Hausdorff? 
        \item Does $Y$ have a countable basis for its 
            topology? 
        \item Check up the definition of connectedness in 
            section 23: is $Y$ connected? 
        \item Check up the definition of compact in section 
            $26$: is $Y$ compact? 
    \end{itemize}
\end{prob*}


\begin{proof}[Solution and proof of all]
    Let
    \[
        q:X\longrightarrow Y=X/{\sim}
    \]
    denote the quotient map. It is useful to write
    \[
        0_-:=q(0,-1),\qquad 0_+:=q(0,1),
    \]
    and, for $a\neq 0$,
    \[
        \overline a:=q(a,-1)=q(a,1).
    \]
    Thus, intuitively, $Y$ is the interval $[-1,1]$ where every
    point except $0$ has been identified between the two copies, while
    the origin occurs twice.

    \medskip

    \textbf{The relation $\sim$ is an equivalence relation.}

    More explicitly, for $(a,s),(b,t)\in X$, we may write
    \[
        (a,s)\sim(b,t)
        \quad\Longleftrightarrow\quad
        a=b\text{ and either }s=t\text{ or }a\neq 0.
    \]

    We check the three conditions.

    \begin{itemize}
        \item \emph{Reflexivity:}
        For every $(a,s)\in X$, clearly
        \[
            (a,s)\sim(a,s).
        \]

        \item \emph{Symmetry:}
        If $(a,s)\sim(b,t)$, then $a=b$. If $s=t$, symmetry is
        immediate. Otherwise $a=b\neq0$, and by the definition of the
        relation we also have
        \[
            (b,t)\sim(a,s).
        \]

        \item \emph{Transitivity:}
        Suppose
        \[
            (a,s)\sim(b,t)
            \qquad\text{and}\qquad
            (b,t)\sim(c,u).
        \]
        Then necessarily $a=b=c$. If this common value is nonzero,
        the two copies are identified, so
        \[
            (a,s)\sim(c,u).
        \]
        If the common value is $0$, the only equivalences at $0$ are
        the reflexive ones. Hence $s=t=u$, and again
        \[
            (a,s)\sim(c,u).
        \]
    \end{itemize}

    Therefore $\sim$ is an equivalence relation.
\end{proof}

\end{document}
